\begin{table}[!htbp]
\centering
\caption{Ratio of direct cross-outcome IV gains to RSS-implied causal gains}
\label{tab:unshrunk-cross-iv-rss-validation-ratios}
\resizebox{\textwidth}{!}{%
\begin{tabular}{lccccccc}
\toprule
Row school VA & Math & Verbal & Exam taking & High-premium inst. & High-premium field & Log income & Fin. aid app. \\
\midrule
Math & 1.00 & 1.23 & 1.81 & 0.54 & -2.28 & -0.65 & 0.22 \\
Verbal & 0.85 & 1.00 & 1.64 & 0.40 & -3.59 & -1.05 & 0.41 \\
Exam taking & 0.76 & 0.54 & 1.00 & -1.07 & -6.36 & -1.13 & 0.24 \\
High-premium inst. & 0.14 & 0.59 & 0.60 & 1.00 & 4.02 & 1.21 & -0.15 \\
High-premium field & 0.06 & 0.76 & -2.27 & 1.48 & 1.00 & 0.80 & 0.50 \\
Log income & 0.32 & 0.87 & -0.14 & 0.97 & 1.26 & 1.00 & 17.76 \\
Fin. aid app. & -0.51 & 0.12 & 1.44 & 0.21 & -5.32 & -9.92 & 1.00 \\
\bottomrule
\end{tabular}}
\par\medskip
\footnotesize
\begin{minipage}{\textwidth}
Notes: Each cell divides the direct unshrunk cross-outcome IV estimate by the corresponding RSS-implied causal prediction, defined as the RSS projection from the row VA to the column VA multiplied by the column outcome's unshrunk same-outcome IV pass-through. A ratio of one denotes exact agreement, zero denotes no direct IV gain despite a nonzero prediction, and a negative ratio denotes opposite signs. Diagonal entries equal one mechanically because their prediction uses the same same-outcome IV coefficient. Ratios can be large when the predicted gain is close to zero; the underlying levels and standard errors are reported in Table~\ref{tab:unshrunk-cross-iv-rss-validation}.
\end{minipage}
\end{table}
